\documentclass[10pt]{article}
\usepackage[margin=1.5cm]{geometry}
\usepackage{amsmath}
\usepackage{graphicx}
\usepackage[spanish]{babel}
\usepackage{parskip}
\pagestyle{empty}

\begin{document}

\begin{minipage}[t][0.23\textheight][t]{0.48\textwidth}
\textbf{Evaluación: Mecánica de los Fluidos} \hfill Nombre: \hrulefill

\vspace{1mm}
\textbf{Curso: 5º Año – Escuela Técnica} \\
\textbf{Duración: 80 minutos}

\vspace{2mm}
\textbf{1.} Calcular la presión absoluta a una profundidad de $h=7,6\,\text{m}$ en agua. ¿Qué presión se ejerce sobre el Titanic a $3843\,\text{m}$?

\vspace{2mm}
\textbf{2.} Determinar la profundidad de agua necesaria para que la presión absoluta sea $2{,}06645\,\text{kg/cm}^2$.

\vspace{2mm}
\textbf{3.} En una prensa hidráulica, el radio del plato menor es $r_1 = 8\,\text{cm}$ y del mayor $r_2 = 50\,\text{cm}$. Si se desea levantar un cuerpo de $1000\,\text{kg}$, ¿qué fuerza hay que aplicar sobre el émbolo menor?

\vspace{2mm}
\textbf{4.} Si $A_1 = 10\,\text{cm}^2$, $A_2 = 100\,\text{cm}^2$ y se aplica una fuerza de $F_2 = 100\,\text{N}$ sobre el émbolo mayor, calcular $F_1$.

\vfill
\textit{Nota: Justificar cada procedimiento. Usar unidades coherentes.}
\end{minipage}
\hspace{0.04\textwidth}
\begin{minipage}[t][0.23\textheight][t]{0.48\textwidth}
% Se puede duplicar para hacer otra versión en la misma hoja
\textbf{Evaluación: Mecánica de los Fluidos} \hfill Nombre: \hrulefill

\vspace{1mm}
\textbf{Curso: 5º Año – Escuela Técnica} \\
\textbf{Duración: 80 minutos}

\vspace{2mm}
\textbf{1.} Calcular la presión absoluta a una profundidad de $h=7,6\,\text{m}$ en agua. ¿Qué presión se ejerce sobre el Titanic a $3843\,\text{m}$?

\vspace{2mm}
\textbf{2.} Determinar la profundidad de agua necesaria para que la presión absoluta sea $2{,}06645\,\text{kg/cm}^2$.

\vspace{2mm}
\textbf{3.} En una prensa hidráulica, el radio del plato menor es $r_1 = 8\,\text{cm}$ y del mayor $r_2 = 50\,\text{cm}$. Si se desea levantar un cuerpo de $1000\,\text{kg}$, ¿qué fuerza hay que aplicar sobre el émbolo menor?

\vspace{2mm}
\textbf{4.} Si $A_1 = 10\,\text{cm}^2$, $A_2 = 100\,\text{cm}^2$ y se aplica una fuerza de $F_2 = 100\,\text{N}$ sobre el émbolo mayor, calcular $F_1$.

\vfill
\textit{Nota: Justificar cada procedimiento. Usar unidades coherentes.}
\end{minipage}

\end{document}
